\begin{frame}[t]{Research record: genomics}
\vspace{0pt}
{\fontsize{8.1}{9.8}\selectfont
\textbf{Genome architecture and bioinformatics}\\
Harris, \dots, \textbf{Eliaz}, \dots, Lieberman Aiden, Rowley. Chromatin alternates between A and B compartments at kilobase scale for subgenic organization. \emph{Nat.\ Commun.} 14:3303 (2023).\\
Kaushal, \dots, \textbf{Eliaz}, \dots. CTCF loss has limited effects on global genome architecture in \emph{Drosophila} despite critical regulatory functions. \emph{Nat.\ Commun.} 12:1011 (2021).\\
Barkai, \dots, \textbf{Eliaz}, \dots. OffRisk: a docker image for annotating CRISPR off-target sites in the human genome. \emph{Bioinformatics Advances} 3:vbad138 (2023).\\[4pt]
\textbf{Cancer genomics}\\
Saurty-Seerunghen, El-Habr, \textbf{Eliaz}, \dots, Junier. A unique malignant cell type per patient tumor encoded in each cancer cell transcriptome. \emph{iScience} 29:115139 (2026).\\
Gilad, \textbf{Eliaz}, \dots. A genome-scale CRISPR Cas9 dropout screen identifies synthetically lethal targets in SRC-3 inhibited cancer cells. \emph{Commun.\ Biol.} 4:399 (2021).\\
Gilad, \textbf{Eliaz}, \dots. Reply to: Target expression is a relevant factor in synthetic lethal screens. \emph{Commun.\ Biol.} 5:836 (2022).\\
Gilad, \textbf{Eliaz}, \dots. Drug-induced PD-L1 expression and cell stress response in breast cancer cells can be balanced by drug combination. \emph{Sci.\ Rep.} 9:15099 (2019).\par}
\sources{Full author lists and positions: CLAIM-FENCE.md.}
\note{[backup] Genome architecture, bioinformatics and cancer genomics at Baylor, NRGene, HIT and Harvard. OffRisk appears on slide 14 (the off-target analogy). The 2021 CRISPR screen appears on slides 33, 40 and 45 (replicates, synthetic lethality). The per-patient tumour clustering (iScience) appears on slides 7 and 37 (cluster labels). The ENCODE pipelines, which set the replay bar on slide 18, are on the next page. Team results with co-authors.}
\end{frame}
